%=====================================================================
\chapter{Pre-registration, Registered Reports and Questionable Practices}\label{ch:prereg}
%=====================================================================
\locator{5}{Part~5 --- Rigour, Reproducibility and Ethics}

\begin{outcomes}
By the end of this chapter you will be able to:
\begin{tight}
  \item \textbf{Identify} the questionable research practices, and explain why
    they are not fraud yet corrupt a literature.
  \item \textbf{Write} a pre-registration that constrains you usefully.
  \item \textbf{Separate} confirmatory from exploratory analysis in a report.
  \item \textbf{Report} deviations from a plan without undermining it.
  \item \textbf{Judge} whether a registered report is available and appropriate
    for your study.
\end{tight}
\end{outcomes}

\begin{prereq}
\chref{inference} for the garden of forking paths and multiplicity;
\chref{sampling} for why the sample size must be fixed in advance;
\chref{slr} --- a review protocol is a pre-registration.
\end{prereq}

\begin{keyterms}
questionable research practice $\bullet$ p-hacking $\bullet$ HARKing $\bullet$
optional stopping $\bullet$ selective reporting $\bullet$ researcher degrees of
freedom $\bullet$ pre-registration $\bullet$ registered report $\bullet$
in-principle acceptance $\bullet$ confirmatory and exploratory $\bullet$
deviation $\bullet$ file drawer
\end{keyterms}

\section{The Practices That Are Not Fraud}

Fabrication and falsification are rare and unambiguous. What corrupts
literatures at scale is a set of practices that feel like ordinary diligence
and are almost never done with intent to deceive.

\begin{center}
\small
\begin{longtable}{@{}L{3.0cm}L{4.8cm}L{6.4cm}@{}}
\toprule
\textbf{Practice} & \textbf{What it looks like from inside} & \textbf{What it does} \\
\midrule
\endfirsthead
\toprule
\textbf{Practice} & \textbf{From inside} & \textbf{What it does} \\
\midrule
\endhead
\bottomrule
\endfoot
p-hacking
  & Trying reasonable analytic variants until one is significant
  & Inflates the false positive rate far above 5\%. With a handful of
    innocuous degrees of freedom it can exceed 60\%~\cite{simmons2011false} \\
HARKing
  & Presenting a hypothesis suggested by the data as if it had been predicted
  & Removes the test. A prediction made after seeing the data has no
    predictive content~\cite{kerr1998harking} \\
Optional stopping
  & Collecting more data until the result stabilises --- which feels prudent
  & Guarantees eventual significance under the null if you keep looking \\
Selective reporting
  & Reporting the outcome measure that worked, describing others as
    preliminary
  & Multiplicity uncorrected and invisible to the reader \\
Post-hoc exclusion
  & Removing participants or runs that were clearly anomalous
  & When the rule is chosen after seeing which way it moves the result, it is
    a form of p-hacking \\
Forking paths
  & Only ever running one analysis --- the obvious one
  & Multiplicity still applies if a different dataset would have prompted a
    different choice~\cite{gelman2013garden} \\
\end{longtable}
\end{center}

\begin{alertbox}[Why good researchers do all of these]
Every one of these has an innocent version. Excluding a participant who
misunderstood the task is correct. Trying a log transform when residuals are
skewed is correct. Collecting more data when the first batch was too small is
sensible. Reporting the analysis that best fits the data's actual structure is
good practice.

\medskip
The corruption is not in the decision; it is in the fact that the decision was
made \emph{after seeing which way it moved the result}, and that the reader is
never told the decision was available. The literature then fills with true-looking
findings that will not replicate, produced entirely by people acting in good
faith.

\medskip
This is why the remedy is procedural rather than moral. You cannot fix it by
resolving to be careful. You fix it by writing the decisions down before you
can be influenced by them.
\end{alertbox}

\section{Pre-registration}

\begin{definitionbox}[Pre-registration]
A time-stamped, publicly deposited statement of your hypotheses, design,
sample size and analysis plan, made before the data are collected --- or, where
data already exist, before you look at them.

\medskip
It does not commit you to a good study, and it does not prevent you from
changing your mind. It makes the change \emph{visible}, which is all that is
required to restore the reader's ability to judge the evidence.
\end{definitionbox}

\begin{center}
\small
\begin{tabular}{@{}L{4.2cm}L{10.1cm}@{}}
\toprule
\textbf{What it fixes} & \textbf{What it does not} \\
\midrule
Distinguishes prediction from postdiction
  & Make a bad design good \\
Fixes the sample size, removing optional stopping
  & Prevent you from exploring --- it just labels the exploration \\
Fixes the primary outcome, removing selective reporting
  & Guarantee the study will be published \\
Makes the analysis decisions visible
  & Apply straightforwardly to all research traditions \\
\bottomrule
\end{tabular}
\end{center}

\begin{templatebox}[C12 --- Pre-registration skeleton]
Deposit on OSF or AsPredicted before collecting. Fifteen minutes of thought per
section is enough; the value is in having decided, not in the prose.

\medskip
\textbf{1. Hypotheses.} Numbered, directional where theory predicts a
direction. Distinguish confirmatory from exploratory now.
\filllines{3}

\noindent\textbf{2. Design.} Conditions, assignment, unit of randomisation,
blinding.
\filllines{2}

\noindent\textbf{3. Sample size and stopping rule.} The number, and the
a-priori power analysis justifying it (\chref{sampling}). ``We will stop at
$N=120$'' --- not ``when the effect stabilises''.
\filllines{2}

\noindent\textbf{4. Variables.} Each construct, its operationalisation, and its
instrument.
\filllines{2}

\noindent\textbf{5. Exclusion rules.} Participant and observation level,
decided now, applied blind to condition.
\filllines{2}

\noindent\textbf{6. Analysis plan.} The exact test or model per hypothesis;
assumption checks and what you will do if they fail; the multiple-comparison
correction and which comparisons it covers.
\filllines{3}

\noindent\textbf{7. Inference criteria.} What result would support the
hypothesis, what would contradict it, and what would be inconclusive.
\filllines{2}

\noindent\textbf{8. Exploratory analyses anticipated.} Listed, and labelled as
exploratory in advance.
\filllines{2}
\end{templatebox}

\begin{conceptbox}[Section 6 is where the value is]
Most people find sections 1--5 easy and section 6 revealing. Writing ``we will
fit a mixed-effects model with reviewer as a random intercept, and if residuals
are non-normal we will report bootstrap intervals alongside'' forces you to
discover, while it is still cheap, that you do not yet know how you will
analyse the data.

\medskip
Writing this section is also the fastest way to find design flaws. If you cannot
specify the analysis, the design is not finished --- which is the same lesson
as the alignment table (\chref{questions}) and the pilot-data analysis script
(\chref{experiments}), arriving for the third time because it is the single
most useful habit in the course.
\end{conceptbox}

\section{Deviations}

Plans meet reality. Deviating is normal; concealing the deviation is not.

\begin{center}
\small
\begin{tabular}{@{}L{4.0cm}L{10.3cm}@{}}
\toprule
\textbf{Do} & \textbf{Do not} \\
\midrule
Report every deviation in a table: planned, actual, reason
  & Silently substitute the analysis that worked \\
State whether the deviation was prompted by the data
  & Describe a data-prompted change as a clarification \\
Report both analyses where feasible --- planned and revised
  & Report only the revised one \\
Relabel affected results as exploratory
  & Keep calling them confirmatory \\
\bottomrule
\end{tabular}
\end{center}

\begin{conceptbox}[The distinction that decides how to report a deviation]
Ask whether the deviation was prompted by something you saw \emph{in the
outcome data}.

\medskip
Discovering that a logging bug corrupted a variable, and dropping it, is not
prompted by the outcome --- report it and continue as confirmatory. Discovering
that the planned test gives $p = 0.07$ and switching to one that gives $p =
0.04$ is prompted by the outcome, and the result is now exploratory whatever
the justification for the new test.

\medskip
Reporting a result as exploratory is not a confession of failure. It is an
accurate label, and readers who know the difference will trust the rest of the
paper more for seeing it.
\end{conceptbox}

\section{Registered Reports}

A stronger mechanism, and one that addresses publication bias directly.

\begin{center}
\small
\begin{tabular}{@{}L{2.2cm}L{12.1cm}@{}}
\toprule
Stage 1 & You submit the introduction, methods and analysis plan ---
          \emph{before} collecting data. Reviewers assess the question's
          importance and the design's adequacy \\
IPA     & On acceptance, the journal issues \emph{in-principle acceptance}: it
          commits to publishing the results whatever they show, provided you
          follow the plan \\
Stage 2 & You collect, analyse and report. Review checks adherence to the plan,
          not whether the findings are exciting \\
\bottomrule
\end{tabular}
\end{center}

\begin{conceptbox}[What in-principle acceptance changes]
It removes the incentive that drives most questionable practice. If publication
is already secured, there is nothing to gain from a significant result, and
p-hacking becomes pointless rather than merely wrong.

\medskip
It also solves the file drawer at source: a null result is published because
the commitment was made before anyone knew it would be null. Evidence from
psychology, where the format is established, shows a far higher proportion of
null results among registered reports than among conventional
papers~\cite{chambers2022past} --- which is what an unbiased literature should
look like.

\medskip
Availability in computing is still limited but growing; some empirical software
engineering venues now run the track. Check whether yours does, because for a
PhD scholar the guarantee of publication regardless of outcome is worth a great
deal against \reg{PhD Regs 2024}{\S15}.
\end{conceptbox}

\section{Does This Apply to All Research?}

Honestly: no, not in the same form, and claiming otherwise damages the case.

\begin{center}
\small
\begin{tabular}{@{}L{3.2cm}L{11.1cm}@{}}
\toprule
Confirmatory quantitative
  & Straightforwardly applicable. Pre-register \\
Systematic review
  & Applicable, and standard practice: the protocol is a pre-registration
    (\chref{slr}). PROSPERO registers reviews in some fields \\
Exploratory quantitative
  & Register the \emph{plan to explore}: what data, what analyses, and that no
    confirmatory claims will be made \\
Qualitative
  & Contested. Emergent design is a methodological commitment, not sloppiness,
    and fixing the analysis in advance would defeat theoretical sampling
    (\chref{qualitative}). Register the research questions, sampling strategy
    and stopping rule; the audit trail (\chref{qda}) does the rest \\
Design science
  & Partially: the objectives and evaluation criteria can and should be fixed
    before building (\chref{designscience}). Iteration on the artefact is the
    method, not a deviation \\
\bottomrule
\end{tabular}
\end{center}

\begin{pitfall}[Pre-registration theatre]
A pre-registration so vague it constrains nothing --- ``we will analyse the data
using appropriate statistical methods'' --- provides the appearance of rigour
with none of the substance, and is arguably worse than none, because it
misleads.

\medskip
The test: could a reader, holding your pre-registration and your dataset,
predict which numbers your paper will report? If not, it was not specific
enough to have prevented anything.
\end{pitfall}

\begin{traditions}{pre-specification}
\tradEMP{Its natural home. Pre-registration converts an analysis into a test
and protects the null result.}
\tradDSR{Fix objectives and evaluation criteria before building; iterate the
artefact freely. Evaluating against criteria invented after the artefact exists
is the failure this prevents.}
\tradQUAL{Resists pre-specified analysis for principled reasons. Registers
questions, sampling logic and stopping rule; relies on the audit trail for the
rest.}
\tradFORM{Not applicable to proofs. Applies to any accompanying empirical
evaluation, which is often the weakest part of a formal paper.}
\end{traditions}

\begin{lab}[Pre-register your study]
\begin{tightnum}
  \item Complete template C12 for your own design.
  \item Section 6 must be specific enough that a colleague, given your raw data
    file, could run the analysis and get your numbers without asking you
    anything.
  \item Give it to a colleague with a simulated dataset of the right shape and
    ask them to do exactly that.
  \item Fix every ambiguity they hit.
  \item Deposit on OSF with a time stamp, and cite it in your synopsis.
\end{tightnum}

\medskip
This is portfolio piece P6. Step 3 is where the ambiguities you cannot see
yourself become visible.
\end{lab}

\begin{checkpoint}
\begin{tightnum}
  \item Why is HARKing damaging even when the hypothesis turns out to be
    correct?
  \item Optional stopping feels like collecting more evidence. Why is it not?
  \item Your planned test's assumptions fail and you switch to a rank-based
    test. Is this a deviation requiring relabelling? What determines the answer?
  \item Why does in-principle acceptance remove the incentive for p-hacking?
\end{tightnum}
\end{checkpoint}

\begin{chaptersummary}
\begin{tight}
  \item Questionable practices are not fraud. They are ordinary decisions made
    after seeing the data, and they corrupt literatures at scale.
  \item The remedy is procedural, not moral: write the decisions down before
    they can be influenced.
  \item Pre-registration fixes hypotheses, sample size, outcomes, exclusions and
    analysis. Section 6 --- the analysis plan --- is where the value is.
  \item Deviations are normal. Report them in a table with reasons, and relabel
    as exploratory anything prompted by the outcome data.
  \item Registered reports secure publication before results exist, removing the
    incentive and the file drawer together.
  \item The mechanism adapts rather than transfers: qualitative work registers
    questions and stopping rules, design science registers objectives and
    evaluation criteria.
  \item A pre-registration that could not have prevented anything is theatre.
\end{tight}
\end{chaptersummary}

\begin{reviewq}
  \item Simmons and colleagues showed that a few innocuous researcher degrees
    of freedom push the false positive rate above 60\%. Enumerate the degrees
    of freedom available in a typical software engineering experiment and
    estimate which contribute most.
  \item Pre-registration is standard in clinical trials and rare in computing.
    Identify the structural differences between the fields that explain this,
    and say which are genuine obstacles and which are merely habit.
  \item Some argue pre-registration privileges confirmatory work and devalues
    discovery. Assess the argument, and say whether the exploratory label
    adequately answers it.
  \item Design a registered report track for a software engineering conference.
    Specify the review timeline, what Stage 1 reviewers assess, and how you
    would handle a study whose artefact turns out not to work.
\end{reviewq}
